\input{preamble}

\title{Section 4.5 --- Exponential Models --- Answer Key}

\begin{document}
\maketitle
\noindent In these exercises, when no rounding instructions are given, round in the way that makes the most sense for the context of the problem.
\section*{Exponential Applications}
\begin{senum}
\item \begin{clist}
\item No: their profits are for a full year.
\item $P(t)=17\cdot 1.2^t$
\item \$35.25 billion
\item 23 years
\end{clist}
%%%%%%%%%%%%%%%%%%%%%%%%%%%%%%%%%%%%%%%%%
\item 
\begin{clist}
\item Yes: cards are destroyed continually throughout the year.
\item $n(t)=10000\cdot 0.88^t$
\item 5277 cards
\item About 72 years
\end{clist}
%%%%%%%%%%%%%%%%%%%%%%%%%%%%%%%%%%%%%%%%%
\item 
\begin{clist}
\item Yes: the Carbon-11 decays continuously in time.
\item $m(t)=27\Gp{\frac{1}{2}}^{t/20}$
\item 3.38 g
\item 142 minutes
\end{clist}
\item A drug is metabolized with the half-life of 7.3 hours. An initial dose of 200mg is initially administered.
\begin{clist}
\item $m(t)=200\Gp{\frac{1}{2}}^{t/7.3}$
\item 31 hours
\item 189 mg
\item 179 mg
\end{clist}
\item The population of a small town declines exponentially after its primary industry goes out of business. In 1997, the population was 30,850. By 2015, it had fallen to only 9,810. (Hint: in setting up this problem, let $t$ represent the number of years since 1997.)
\begin{clist}
\item $p(t)=30850\Gp{\dfrac{9810}{30850}}^{t/18}$
\item 17,397 people
\item 2003
\item 6.2\%
\end{clist}
\end{senum}
\section*{Compound Interest}
You may wish to reference the compound interest formula in solving these problems:
$$P(t)=P_0\left(1+\frac{r}{n}\right)^{nt}$$
\begin{smenum}
\mitemxxxx
{\$2902.01}
{\$2817.48}
{\$1897.98}
{\$779.71}
\end{smenum}
\begin{senum}
%%%%%%%%%%%%%%%%%%%%%%%%%%%%%%%%%%%%%%%%%
\item
\begin{nmenum}
\mitemxxx
{\$22,921.83}
{23.191 years}
{By 5.776 years}
\mitemx
{36.095 years at 3\% interest, 27.105 years at 4\% interest}
\end{nmenum}
%%%%%%%%%%%%%%%%%%%%%%%%%%%%%%%%%%%%%%%%%
\item 
\begin{nmenum}
\mitemxxo
{\$2199.76}
{180 years}
\end{nmenum}
\end{senum}

\end{document}